\documentclass[12pt,a4paper]{article}
\usepackage[UTF8]{ctex}
\usepackage{amsmath,amssymb,amsthm}
\usepackage{geometry}
\usepackage{graphicx}
\usepackage{hyperref}
\usepackage{bm}

\geometry{margin=2.5cm}

\title{关节力矩公式$\tau_i = F_i^T A_i$的推导}
\author{Modern Robotics}
\date{}

\begin{document}

\maketitle

\section{引言}

在牛顿-欧拉逆向动力学算法中，关节力矩的计算公式为：
\begin{equation}
\tau_i = F_i^T A_i
\label{eq:main_formula}
\end{equation}

其中：
\begin{itemize}
\item $\tau_i \in \mathbb{R}$：关节$i$的力矩（旋转关节）或力（移动关节）
\item $F_i \in \mathbb{R}^6$：通过关节$i$传递到连杆$i$的wrench（力和力矩的组合）
\item $A_i \in \mathbb{R}^6$：关节$i$的螺旋轴（screw axis）
\end{itemize}

本文档将详细推导这个公式的来源，从物理原理和数学投影两个角度进行解释。

\section{基本概念}

\subsection{Wrench的定义}

\textbf{Wrench}（力螺旋）是力和力矩的组合，用6维向量表示：
\begin{equation}
F_i = \begin{bmatrix}
m_i \\
f_i
\end{bmatrix} \in \mathbb{R}^6
\end{equation}

其中：
\begin{itemize}
\item $m_i \in \mathbb{R}^3$：作用在连杆$i$上的力矩
\item $f_i \in \mathbb{R}^3$：作用在连杆$i$上的力
\end{itemize}

在牛顿-欧拉算法中，$F_i$表示通过关节$i$传递到连杆$i$的总wrench，包括：
\begin{enumerate}
\item 连杆$i$自身的惯性力
\item 从连杆$i+1$传递过来的力（通过关节$i+1$）
\item 末端执行器施加的外力（对于末端连杆）
\end{enumerate}

\subsection{螺旋轴（Screw Axis）的定义}

\textbf{螺旋轴} $A_i$描述了关节$i$的运动方向和类型，也是一个6维向量：
\begin{equation}
A_i = \begin{bmatrix}
\omega_{A_i} \\
v_{A_i}
\end{bmatrix} \in \mathbb{R}^6
\end{equation}

其中：
\begin{itemize}
\item \textbf{旋转关节}：$\omega_{A_i} \in \mathbb{R}^3$是单位旋转轴方向，$v_{A_i} = 0$
\item \textbf{移动关节}：$\omega_{A_i} = 0$，$v_{A_i} \in \mathbb{R}^3$是单位移动方向
\end{itemize}

螺旋轴$A_i$通常被归一化为单位向量，即$\|A_i\| = 1$。

\section{物理原理：关节约束}

\subsection{关节的运动约束}

关节$i$的物理特性决定了它\textbf{只能沿螺旋轴$A_i$方向传递力或力矩}。具体来说：

\begin{enumerate}
\item \textbf{旋转关节}：
\begin{itemize}
\item 只能传递沿旋转轴方向的力矩
\item 不能传递垂直于旋转轴的力矩（被轴承约束）
\item 不能传递任何力（假设关节是理想的，无摩擦）
\end{itemize}

\item \textbf{移动关节}：
\begin{itemize}
\item 只能传递沿移动方向的力
\item 不能传递垂直于移动方向的力（被导轨约束）
\item 不能传递任何力矩（假设关节是理想的）
\end{itemize}
\end{enumerate}

\subsection{关节约束的数学表达}

由于关节只能沿$A_i$方向传递力/力矩，wrench $F_i$中只有\textbf{沿$A_i$方向的分量}对关节运动有贡献。

设$F_i$可以分解为：
\begin{equation}
F_i = F_{i,\parallel} + F_{i,\perp}
\end{equation}

其中：
\begin{itemize}
\item $F_{i,\parallel}$：沿$A_i$方向的分量（可以被关节传递）
\item $F_{i,\perp}$：垂直于$A_i$方向的分量（被关节约束，无法传递）
\end{itemize}

关节力矩$\tau_i$就是$F_{i,\parallel}$在$A_i$方向上的大小。

\subsection{重要澄清：为什么是"沿$A_i$方向"？}

\textbf{常见误解}：有人认为应该是$\bm{r} \times \bm{F}$（力臂叉乘力）产生力矩，为什么这里说是"沿$A_i$方向的分量"？

\textbf{关键理解}：

\begin{enumerate}
\item \textbf{Wrench $F_i$已经包含了所有力矩信息}

$F_i = \begin{bmatrix} m_i \\ f_i \end{bmatrix}$中的$m_i$是作用在连杆$i$上的\textbf{总力矩}，这个力矩可能来自：
\begin{itemize}
\item 直接作用的力矩
\item 力$\bm{f}$通过力臂$\bm{r}$产生的力矩：$\bm{r} \times \bm{f}$
\item 其他连杆传递过来的力矩
\item 惯性力矩等
\end{itemize}

在计算$F_i$时（通过牛顿-欧拉后向递推），所有这些力矩贡献已经被汇总到$m_i$中了。

\item \textbf{对于旋转关节：只有力矩在旋转轴方向的分量能被传递}

对于旋转关节，$A_i = \begin{bmatrix} \omega_{A_i} \\ 0 \end{bmatrix}$，因此：
\begin{align}
\tau_i &= F_i^T A_i = \begin{bmatrix} m_i^T & f_i^T \end{bmatrix} \begin{bmatrix} \omega_{A_i} \\ 0 \end{bmatrix} \\
&= m_i^T \omega_{A_i} = \bm{m}_i \cdot \hat{\bm{\omega}}_i
\end{align}

这表示：\textbf{只有力矩向量$\bm{m}_i$在旋转轴$\hat{\bm{\omega}}_i$方向上的分量能被关节传递}。

\item \textbf{物理原因}

旋转关节的物理结构（轴承）决定了：
\begin{itemize}
\item \textbf{可以传递}：沿旋转轴方向的力矩（使关节旋转）
\item \textbf{不能传递}：垂直于旋转轴的力矩（被轴承约束，产生反作用力）
\item \textbf{不能传递}：任何力（假设理想关节，无摩擦）
\end{itemize}

\item \textbf{与$\bm{r} \times \bm{F}$的关系}

$\bm{r} \times \bm{f}$确实会产生力矩，但这个力矩的贡献已经包含在$m_i$中了。关键点是：
\begin{itemize}
\item $\bm{r} \times \bm{f}$产生的力矩是一个3维向量
\item 这个力矩向量可能不在旋转轴方向上
\item 只有这个力矩向量在旋转轴方向上的分量（即$(\bm{r} \times \bm{f}) \cdot \hat{\bm{\omega}}_i$）才能被关节传递
\item 垂直于旋转轴的力矩分量被轴承吸收，无法传递
\end{itemize}

\item \textbf{6维空间中的投影}

"沿$A_i$方向"指的是在\textbf{6维wrench空间}中的投影：
\begin{equation}
F_{i,\parallel} = (F_i^T A_i) A_i
\end{equation}

对于旋转关节，这等价于：
\begin{equation}
F_{i,\parallel} = \begin{bmatrix} (m_i^T \omega_{A_i}) \omega_{A_i} \\ 0 \end{bmatrix} = \begin{bmatrix} (\bm{m}_i \cdot \hat{\bm{\omega}}_i) \hat{\bm{\omega}}_i \\ 0 \end{bmatrix}
\end{equation}

即：力矩部分只保留在旋转轴方向上的分量，力部分全部被约束（理想关节）。
\end{enumerate}

\textbf{总结}：
\begin{itemize}
\item $F_i$中的$m_i$已经包含了所有力矩贡献（包括$\bm{r} \times \bm{f}$）
\item 对于旋转关节，只有$m_i$在旋转轴$\omega_{A_i}$方向上的分量能被传递
\item "沿$A_i$方向"指的是在6维wrench空间中沿螺旋轴方向的投影
\item 这与$\bm{r} \times \bm{F}$不矛盾，而是$\bm{r} \times \bm{F}$的贡献已经包含在$m_i$中，然后我们提取$m_i$在旋转轴方向的分量
\end{itemize}

\subsubsection{具体例子：理解"沿$A_i$方向"的含义}

\textbf{场景}：考虑一个旋转关节，旋转轴沿$z$轴方向，即$\hat{\bm{\omega}}_i = \begin{bmatrix} 0 \\ 0 \\ 1 \end{bmatrix}$。

\textbf{情况1}：假设wrench $F_i$为：
\begin{equation}
F_i = \begin{bmatrix} m_i \\ f_i \end{bmatrix} = \begin{bmatrix} 2 \\ 1 \\ 3 \\ 0 \\ 0 \\ 0 \end{bmatrix}
\end{equation}

其中力矩$\bm{m}_i = \begin{bmatrix} 2 \\ 1 \\ 3 \end{bmatrix}$，力$\bm{f}_i = \begin{bmatrix} 0 \\ 0 \\ 0 \end{bmatrix}$。

\textbf{计算关节力矩}：
\begin{align}
\tau_i &= F_i^T A_i = \begin{bmatrix} 2 & 1 & 3 & 0 & 0 & 0 \end{bmatrix} \begin{bmatrix} 0 \\ 0 \\ 1 \\ 0 \\ 0 \\ 0 \end{bmatrix} = 3
\end{align}

\textbf{解释}：
\begin{itemize}
\item 力矩向量$\bm{m}_i = \begin{bmatrix} 2 \\ 1 \\ 3 \end{bmatrix}$有三个分量
\item 只有$z$方向的分量（值为$3$）能被关节传递
\item $x$和$y$方向的力矩分量（$2$和$1$）被轴承约束，无法传递
\end{itemize}

\textbf{情况2}：假设力矩来自力通过力臂产生

假设有一个力$\bm{f} = \begin{bmatrix} 1 \\ 0 \\ 0 \end{bmatrix}$作用在点$\bm{r} = \begin{bmatrix} 0 \\ 0 \\ 2 \end{bmatrix}$处（相对于关节中心）。

\textbf{计算力矩}：
\begin{align}
\bm{m} = \bm{r} \times \bm{f} &= \begin{bmatrix} 0 \\ 0 \\ 2 \end{bmatrix} \times \begin{bmatrix} 1 \\ 0 \\ 0 \end{bmatrix} \\
&= \begin{bmatrix} 0 \cdot 0 - 2 \cdot 0 \\ 2 \cdot 1 - 0 \cdot 0 \\ 0 \cdot 0 - 0 \cdot 1 \end{bmatrix} = \begin{bmatrix} 0 \\ 2 \\ 0 \end{bmatrix}
\end{align}

这个力矩$\bm{m} = \begin{bmatrix} 0 \\ 2 \\ 0 \end{bmatrix}$在$y$方向上，\textbf{不在旋转轴$z$方向上}。

\textbf{如果这是wrench中的力矩}：
\begin{equation}
F_i = \begin{bmatrix} 0 \\ 2 \\ 0 \\ 1 \\ 0 \\ 0 \end{bmatrix}
\end{equation}

\textbf{计算关节力矩}：
\begin{align}
\tau_i &= F_i^T A_i = \begin{bmatrix} 0 & 2 & 0 & 1 & 0 & 0 \end{bmatrix} \begin{bmatrix} 0 \\ 0 \\ 1 \\ 0 \\ 0 \\ 0 \end{bmatrix} = 0
\end{align}

\textbf{结果}：$\tau_i = 0$！

\textbf{解释}：
\begin{itemize}
\item $\bm{r} \times \bm{f}$确实产生了力矩$\begin{bmatrix} 0 \\ 2 \\ 0 \end{bmatrix}$
\item 但这个力矩在$y$方向上，\textbf{垂直于旋转轴$z$}
\item 因此，这个力矩\textbf{无法被关节传递}（被轴承约束）
\item 关节力矩$\tau_i = 0$，因为力矩在旋转轴方向上的分量为$0$
\end{itemize}

\textbf{关键理解}：
\begin{enumerate}
\item $\bm{r} \times \bm{f}$产生的力矩已经包含在$m_i$中
\item 但只有$m_i$在旋转轴方向上的分量能被传递
\item 如果$\bm{r} \times \bm{f}$产生的力矩垂直于旋转轴，它对关节力矩的贡献为$0$
\item "沿$A_i$方向"指的是力矩向量$\bm{m}_i$在旋转轴$\hat{\bm{\omega}}_i$方向上的投影
\end{enumerate}

\section{数学推导：向量投影}

\subsection{向量投影公式}

给定向量$\bm{F}$和单位向量$\bm{A}$（$\|\bm{A}\| = 1$），$\bm{F}$在$\bm{A}$方向上的投影为：
\begin{equation}
\text{proj}_{\bm{A}}(\bm{F}) = (\bm{F}^T \bm{A}) \bm{A}
\end{equation}

投影的\textbf{标量大小}为：
\begin{equation}
\|\text{proj}_{\bm{A}}(\bm{F})\| = |\bm{F}^T \bm{A}|
\end{equation}

如果$\bm{A}$是单位向量，则投影的标量大小就是$\bm{F}^T \bm{A}$。

\subsection{应用到关节力矩计算}

对于关节$i$，我们需要计算wrench $F_i$在螺旋轴$A_i$方向上的投影。

假设$A_i$是单位向量（$\|A_i\| = 1$），则$F_i$在$A_i$方向上的投影大小为：
\begin{equation}
F_{i,\parallel} = (F_i^T A_i) A_i
\end{equation}

投影的标量大小为：
\begin{equation}
|F_i^T A_i|
\end{equation}

由于$A_i$是单位向量，这个标量值就是$F_i^T A_i$。

\subsection{符号约定}

关节力矩$\tau_i$的符号约定：
\begin{itemize}
\item 如果$F_i^T A_i > 0$，力矩/力的方向与$A_i$方向相同
\item 如果$F_i^T A_i < 0$，力矩/力的方向与$A_i$方向相反
\end{itemize}

因此，关节力矩为：
\begin{equation}
\tau_i = F_i^T A_i
\label{eq:torque_projection}
\end{equation}

\section{虚功原理推导}

\subsection{虚功原理}

\textbf{虚功原理}（Principle of Virtual Work）是分析力学中的基本原理，它指出：在平衡状态下，所有约束力在虚位移上做的虚功为零。

对于关节$i$，考虑一个虚位移（或虚速度）$\delta \dot{\theta}_i$，这会产生一个虚twist：
\begin{equation}
\delta V_i = A_i \delta \dot{\theta}_i
\end{equation}

\subsection{功率平衡}

作用在连杆$i$上的wrench $F_i$在虚twist $\delta V_i$上产生的虚功率为：
\begin{equation}
P = F_i^T \delta V_i = F_i^T (A_i \delta \dot{\theta}_i) = (F_i^T A_i) \delta \dot{\theta}_i
\end{equation}

另一方面，关节力矩$\tau_i$在虚速度$\delta \dot{\theta}_i$上产生的虚功率为：
\begin{equation}
P = \tau_i \delta \dot{\theta}_i
\end{equation}

根据虚功原理，这两个功率必须相等（因为关节约束，垂直于$A_i$方向的力不做功）：
\begin{equation}
\tau_i \delta \dot{\theta}_i = (F_i^T A_i) \delta \dot{\theta}_i
\end{equation}

由于$\delta \dot{\theta}_i$是任意的，我们得到：
\begin{equation}
\tau_i = F_i^T A_i
\label{eq:virtual_work}
\end{equation}

\section{具体例子}

\subsection{旋转关节}

对于旋转关节$i$，螺旋轴为：
\begin{equation}
A_i = \begin{bmatrix}
\omega_{A_i} \\
0
\end{bmatrix} = \begin{bmatrix}
\hat{\omega}_i \\
0
\end{bmatrix}
\end{equation}

其中$\hat{\omega}_i$是单位旋转轴方向向量。

wrench为：
\begin{equation}
F_i = \begin{bmatrix}
m_i \\
f_i
\end{bmatrix}
\end{equation}

关节力矩为：
\begin{align}
\tau_i &= F_i^T A_i \\
&= \begin{bmatrix} m_i^T & f_i^T \end{bmatrix} \begin{bmatrix} \hat{\omega}_i \\ 0 \end{bmatrix} \\
&= m_i^T \hat{\omega}_i \\
&= \bm{m}_i \cdot \hat{\bm{\omega}}_i
\end{align}

这表示关节力矩是力矩向量$m_i$在旋转轴方向$\hat{\omega}_i$上的投影。

\textbf{物理意义}：
\begin{itemize}
\item 只有沿旋转轴方向的力矩分量对关节运动有贡献
\item 垂直于旋转轴的力矩分量被轴承约束，无法传递
\item 力$f_i$对旋转关节的力矩没有贡献（假设关节是理想的）
\end{itemize}

\subsection{移动关节}

对于移动关节$i$，螺旋轴为：
\begin{equation}
A_i = \begin{bmatrix}
0 \\
v_{A_i}
\end{bmatrix} = \begin{bmatrix}
0 \\
\hat{v}_i
\end{bmatrix}
\end{equation}

其中$\hat{v}_i$是单位移动方向向量。

wrench为：
\begin{equation}
F_i = \begin{bmatrix}
m_i \\
f_i
\end{bmatrix}
\end{equation}

关节力为：
\begin{align}
\tau_i &= F_i^T A_i \\
&= \begin{bmatrix} m_i^T & f_i^T \end{bmatrix} \begin{bmatrix} 0 \\ \hat{v}_i \end{bmatrix} \\
&= f_i^T \hat{v}_i \\
&= \bm{f}_i \cdot \hat{\bm{v}}_i
\end{align}

这表示关节力是力向量$f_i$在移动方向$\hat{v}_i$上的投影。

\textbf{物理意义}：
\begin{itemize}
\item 只有沿移动方向的力分量对关节运动有贡献
\item 垂直于移动方向的力分量被导轨约束，无法传递
\item 力矩$m_i$对移动关节的力没有贡献（假设关节是理想的）
\end{itemize}

\section{在牛顿-欧拉算法中的应用}

\subsection{后向递推中的使用}

在牛顿-欧拉逆向动力学算法的后向递推阶段：

\begin{enumerate}
\item 首先计算wrench $F_i$：
\begin{equation}
F_i = \text{Ad}_{T_{i+1,i}}^T(F_{i+1}) + G_i \dot{V}_i - [\text{ad}_{V_i}]^T G_i V_i
\label{eq:wrench_calculation}
\end{equation}

\item 然后计算关节力矩：
\begin{equation}
\tau_i = F_i^T A_i
\label{eq:torque_calculation}
\end{equation}
\end{enumerate}

\subsection{计算流程}

完整的计算流程为：

\begin{enumerate}
\item \textbf{前向递推}：从基座到末端，计算每个连杆的速度$V_i$和加速度$\dot{V}_i$

\item \textbf{后向递推}：从末端到基座
\begin{enumerate}
\item 对于末端连杆$n$：$F_{n+1} = F_{\text{tip}}$（末端外力，如果没有则为0）
\item 对于每个连杆$i = n, n-1, \ldots, 1$：
\begin{itemize}
\item 计算$F_i$（使用式\ref{eq:wrench_calculation}）
\item 计算$\tau_i = F_i^T A_i$（使用式\ref{eq:torque_calculation}）
\end{itemize}
\end{enumerate}
\end{enumerate}

\section{几何解释}

\subsection{内积的几何意义}

公式$\tau_i = F_i^T A_i$是6维空间中的内积（点积）：
\begin{equation}
F_i^T A_i = \|F_i\| \|A_i\| \cos \theta
\end{equation}

其中$\theta$是$F_i$和$A_i$之间的夹角。

由于$A_i$通常是单位向量（$\|A_i\| = 1$），我们有：
\begin{equation}
\tau_i = \|F_i\| \cos \theta
\end{equation}

这表示$\tau_i$是$F_i$在$A_i$方向上的\textbf{有向投影长度}。

\subsection{投影的分解}

wrench $F_i$可以分解为：
\begin{equation}
F_i = \underbrace{(F_i^T A_i) A_i}_{F_{i,\parallel}} + \underbrace{F_i - (F_i^T A_i) A_i}_{F_{i,\perp}}
\end{equation}

其中：
\begin{itemize}
\item $F_{i,\parallel} = (F_i^T A_i) A_i$：沿$A_i$方向的分量，大小为$\tau_i = F_i^T A_i$
\item $F_{i,\perp} = F_i - (F_i^T A_i) A_i$：垂直于$A_i$方向的分量，被关节约束
\end{itemize}

关节只能传递$F_{i,\parallel}$，而$F_{i,\perp}$被关节的物理结构（轴承、导轨等）吸收。

\section{总结}

关节力矩公式$\tau_i = F_i^T A_i$的来源可以总结为：

\begin{enumerate}
\item \textbf{物理约束}：关节只能沿螺旋轴$A_i$方向传递力/力矩

\item \textbf{数学投影}：$\tau_i$是wrench $F_i$在螺旋轴$A_i$方向上的投影大小

\item \textbf{虚功原理}：通过虚功原理，可以证明$\tau_i \delta \dot{\theta}_i = F_i^T (A_i \delta \dot{\theta}_i)$，从而得到$\tau_i = F_i^T A_i$

\item \textbf{几何意义}：$\tau_i$是6维空间中$F_i$和$A_i$的内积，表示$F_i$在$A_i$方向上的有向投影长度
\end{enumerate}

这个公式是牛顿-欧拉逆向动力学算法的核心，它将6维的wrench $F_i$转换为标量的关节力矩$\tau_i$，使得我们可以计算驱动每个关节所需的力矩或力。

\section{参考文献}

\begin{itemize}
\item Modern Robotics: Mechanics, Planning, and Control, Chapter 8
\item 牛顿-欧拉逆向动力学详细推导：\texttt{newton\_euler.tex}
\end{itemize}

\end{document}

